% Lean: EconLean/BoundedRationality/{LevelK,QRE,BerkNash,Existence}.lean
\chapter{Bounded Rationality in Games}
\label{chap:bounded-rationality}

Sources: Nagel (1995); Stahl--Wilson (1995); Camerer--Ho--Chong (2004); McKelvey--Palfrey (1995);
Esponda--Pouzo (2016).

\section{Level-k and cognitive hierarchy}
\label{sec:br-levelk}

\begin{definition}[Finite game, best response, level-k hierarchy]
\label{def:br-game}
\lean{Econ.BoundedRationality.FiniteGame, Econ.BoundedRationality.BestResponse, Econ.BoundedRationality.IsPureNash, Econ.BoundedRationality.IsStrictNash, Econ.BoundedRationality.IsLevelKPath}
\leanok
Level-0 is a pure anchor profile (not uniform randomization). A level-k hierarchy is a sequence of pure profiles $L_k$ with $L_{k+1}(i)$ a best response of $i$ to $L_k$.
\end{definition}

\begin{theorem}[Level-k basics]
\label{thm:br-levelk}
\lean{Econ.BoundedRationality.exists_bestResponse, Econ.BoundedRationality.exists_levelKPath, Econ.BoundedRationality.isPureNash_of_levelK_fixed, Econ.BoundedRationality.levelK_eq_of_strictNash}
\leanok
\uses{def:br-game}
A hierarchy exists from any level-0 profile; a stationary hierarchy is at a pure Nash equilibrium; if level-0 is
a strict Nash equilibrium every level plays it.
\end{theorem}
\begin{proof}
\leanok
Choose best responses (they exist on finite sets). At a fixed point $L_{k+1}=L_k$ each player's own action is a
best response. Under strictness any other best response would strictly lose, by induction on $k$.
\end{proof}

\begin{theorem}[Beauty contest and cognitive hierarchy]
\label{thm:br-beauty}
\lean{Econ.BoundedRationality.levelK_guess, Econ.BoundedRationality.levelK_guess_tendsto, Econ.BoundedRationality.beauty_fixed_iff, Econ.BoundedRationality.chGuess_strictAnti, Econ.BoundedRationality.chGuess_pos, Econ.BoundedRationality.chGuess_succ, Econ.BoundedRationality.chAvg_succ_lt}
\leanok
In the large-population approximation (best response $px$ to others' mean $x$), level-k guesses are $50p^k\to0$,
the unique fixed point of $x\mapsto px$ (the finite game itself is not formalized). CH guesses ($p$ times the weighted mean
of lower levels, positive weights) are positive and strictly decreasing in the level.
\end{theorem}
\begin{proof}
\leanok
The weighted mean of levels $0..k+1$ averages the mean of $0..k$ with $p$ times that mean, so it falls strictly;
the next guess is $p$ times that mean.
\end{proof}

\section{Quantal response equilibrium}
\label{sec:br-qre}

\begin{definition}[Logit QRE]
\label{def:br-qre}
\lean{Econ.BoundedRationality.U, Econ.BoundedRationality.logitResponse, Econ.BoundedRationality.IsLogitQRE, Econ.BoundedRationality.IsMixedNash}
\leanok
\uses{def:br-game}
$\sigma_i(b)=e^{\lambda U_i(b,\sigma_{-i})}/\sum_ce^{\lambda U_i(c,\sigma_{-i})}$ for all $i,b$.
\end{definition}

\begin{theorem}[McKelvey--Palfrey: limits are Nash]
\label{thm:br-qre-limit}
\lean{Econ.BoundedRationality.isNash_of_tendsto_qre, Econ.BoundedRationality.qre_lambda_zero, Econ.BoundedRationality.uniform_isLogitQRE_zero, Econ.BoundedRationality.sum_exp_pos, Econ.BoundedRationality.IsLogitQRE.isMixedProfile, Econ.BoundedRationality.tendsto_U, Econ.BoundedRationality.logitResponse_le}
\leanok
\uses{def:br-qre}
At $\lambda=0$ the QRE is exactly uniform play. If $\lambda_n\to\infty$ and logit QRE $\sigma_n\to\sigma$, then $\sigma$ is a mixed
Nash equilibrium.
\end{theorem}
\begin{proof}
\leanok
Expected payoffs are continuous. If $\sigma_i(b)>0$ but $U_i(c,\sigma)>U_i(b,\sigma)$, the gap exceeds $\delta/2>0$
eventually, so $\sigma_{n,i}(b)\le e^{-\lambda_n\delta/2}\to0$, contradicting $\sigma_i(b)>0$.
\end{proof}

\begin{theorem}[Existence in 2×2 games]
\label{thm:br-qre-2x2}
\lean{Econ.BoundedRationality.qre2x2_exists, Econ.BoundedRationality.isLogitQRE_of_qre2x2, Econ.BoundedRationality.exists_isLogitQRE_2x2, Econ.BoundedRationality.game2x2, Econ.BoundedRationality.profile2x2, Econ.BoundedRationality.U_row, Econ.BoundedRationality.U_col, Econ.BoundedRationality.exp_div_exp_add_exp, Econ.BoundedRationality.exp_div_exp_add_exp', Econ.BoundedRationality.sum_fin2_arrow, Econ.BoundedRationality.logistic, Econ.BoundedRationality.continuous_logistic, Econ.BoundedRationality.rowAdv, Econ.BoundedRationality.colAdv, Econ.BoundedRationality.IsQRE2x2, Econ.BoundedRationality.logistic_pos, Econ.BoundedRationality.logistic_lt_one}
\leanok
Every $2\times2$ game has a logit QRE for every $\lambda$, in the general sense of Definition~\ref{def:br-qre}.
\end{theorem}
\begin{proof}
\leanok
The composite response $p\mapsto F(p)$ is continuous from $[0,1]$ into $(0,1)$, so $F(p)-p$ changes sign;
apply the intermediate value theorem. With two actions the logit probability of the first is the logistic
function of the payoff difference, and the payoff differences are \texttt{rowAdv}, \texttt{colAdv}.
\end{proof}

\section{Berk--Nash equilibrium}
\label{sec:br-berknash}

\begin{definition}[Berk--Nash equilibrium]
\label{def:br-bn}
\lean{Econ.BoundedRationality.SubjectiveProblem, Econ.BoundedRationality.SubjectiveProblem.K, Econ.BoundedRationality.SubjectiveProblem.IsBerkNash}
\leanok
$(\sigma,\mu)$: actions in the support of $\sigma$ maximize subjective payoff under $\mu$, and $\mu$ is supported on
minimizers of $K(\sigma,\theta)=\sum_x\sigma_x\mathrm{KL}(Q_x\|Q_{\theta,x})$.
\end{definition}

\begin{theorem}[Correct specification]
\label{thm:br-bn-correct}
\lean{Econ.BoundedRationality.SubjectiveProblem.K_nonneg, Econ.BoundedRationality.SubjectiveProblem.K_correct, Econ.BoundedRationality.SubjectiveProblem.isBerkNash_of_correct, Econ.BoundedRationality.SubjectiveProblem.optimal_of_isBerkNash, Econ.BoundedRationality.SubjectiveProblem.pointBelief_isProb, Econ.BoundedRationality.SubjectiveProblem.subjEU_pointBelief}
\leanok
\uses{def:br-bn, thm:beh-mi}
$K\ge0$. If $Q_{\theta^*}=Q$, $\theta^*$ minimizes $K$, optimal play with belief $\delta_{\theta^*}$ is Berk--Nash, and
under identification every Berk--Nash action is optimal.
\end{theorem}
\begin{proof}
\leanok
Gibbs' inequality. Believed parameters have $K\le K(\theta^*)=0$, so by identification they equal the truth,
making subjective and true payoffs equal.
\end{proof}

\begin{theorem}[Misspecification sustains a worse action]
\label{thm:br-bn-example}
\lean{Econ.BoundedRationality.SubjectiveProblem.misspecified_example, Econ.BoundedRationality.SubjectiveProblem.example1_misspecified}
\leanok
\uses{def:br-bn}
Two actions with true success probabilities $1/2$ and $7/10$ (the second costs $1/10$). The model says the action
does not matter and the success probability is $2/5$ or $9/20$, so no parameter is correct. Playing action $0$ with
belief on $9/20$ (the KL-best fit) is Berk--Nash, although action $1$ is truly better.
\end{theorem}
\begin{proof}
\leanok
Under $9/20$ the subjective payoffs are $9/20$ and $9/20-1/10$. The fit comparison reduces to
$\log(2/5)+\log(3/5)\le\log(9/20)+\log(11/20)$, i.e. $0.24\le0.2475$ by \texttt{Real.log\_le\_log}. True payoffs are
$1/2<3/5$.
\end{proof}

\section{Existence}
\label{sec:br-existence}

\begin{theorem}[McKelvey--Palfrey 1995, Theorem 1]
\label{thm:br-qre-exists}
\lean{Econ.BoundedRationality.exists_isLogitQRE, Econ.BoundedRationality.continuous_U, Econ.BoundedRationality.mixedProfiles, Econ.BoundedRationality.convex_mixedProfiles, Econ.BoundedRationality.isCompact_mixedProfiles}
\leanok
\uses{def:br-qre}
Every finite game with nonempty action sets has a logit QRE at every precision $\lambda$.
\end{theorem}
\begin{proof}
\leanok
Mixed profiles form a nonempty compact convex subset of $\prod_i\mathbb{R}^{A_i}$ (closed, inside a product of
unit cubes). The logit response map is continuous (expected payoffs are polynomials in $\sigma$) and lands in mixed
profiles. Brouwer's theorem gives a fixed point.
\end{proof}

\begin{theorem}[Esponda--Pouzo 2016, Theorem 1, finite case]
\label{thm:br-bn-exists}
\lean{Econ.BoundedRationality.SubjectiveProblem.exists_isBerkNash, Econ.BoundedRationality.SubjectiveProblem.probs, Econ.BoundedRationality.SubjectiveProblem.convex_probs, Econ.BoundedRationality.SubjectiveProblem.isClosed_probs, Econ.BoundedRationality.SubjectiveProblem.isCompact_probs, Econ.BoundedRationality.SubjectiveProblem.exists_prob_supported_max, Econ.BoundedRationality.SubjectiveProblem.continuous_K, Econ.BoundedRationality.SubjectiveProblem.continuous_subjEU}
\leanok
\uses{def:br-bn}
With finitely many actions and parameters (and full-support model distributions), a Berk--Nash equilibrium exists.
\end{theorem}
\begin{proof}
\leanok
On $\Delta(X)\times\Delta(\Theta)$ map $(\sigma,\mu)$ to the pairs $(\sigma',\mu')$ with $\sigma'$ supported on
subjective best responses to $\mu$ and $\mu'$ on KL-minimizers given $\sigma$. A point mass on an argmax shows values
are nonempty. Writing ``supported on maximizers of $f$'' as $\sigma'_x(f_{x'}-f_x)\le0$ for all $x,x'$ shows values
are convex and the graph is closed ($K$ and subjective payoffs are continuous). Kakutani's theorem gives a fixed point.
\end{proof}
